% !TeX root = ../main.tex

\chapter{Technische architectuur}
\label{chap:architectuur}

\pending{Dit hoofdstuk inhoudelijk nakijken en afwerken. Verifieer de beschreven modulegrenzen, instellingen en datastroom aan de hand van de definitieve implementatie en werk de verwijzing naar de datastroomfiguur af.}

\section{Electron}
De applicatie is gebouwd met Electron en JavaScript. Electron combineert een Chromium-renderer met
een Node.js-mainproces. Deze keuze maakte het mogelijk om een webgebaseerde interface te bouwen en
tegelijk lokale bestanden en meerdere desktopvensters te beheren. Dezelfde
codebasis kan voor macOS en Windows worden verpakt.

Het mainproces in \code{src/main/main.js} beheert de levenscyclus van de applicatie, instellingen,
vensters, polling, opslag en IPC-handlers. Het verborgen livevenster laadt de echte timingwebsite.
Een extractiescript wordt in die pagina uitgevoerd om tabelkoppen, cellen en sessie-informatie uit
de gerenderde DOM te lezen.

\section{Proces- en modulegrenzen}
De code is bewust in vier zones verdeeld:

\begin{itemize}
  \item \code{src/main}: Electron-specifieke infrastructuur, vensters, updates en bestands-I/O.
  \item \code{src/shared}: pure domeinfuncties voor parsing, analyse, voorspelling, pitstops en
        viewmodellen.
  \item \code{src/renderer}: HTML, CSS en code die reeds berekende data op het scherm zet.
  \item \code{tests}: alle testfiles om domeinfuncties als ook renderer en main functionaliteiten grondig te testen.
\end{itemize}

De scheiding tussen renderer en domeinlogica was een belangrijke verbetering tijdens de stage.
Vroege versies bevatten berekeningen in de UI-code. Daardoor was het moeilijk om te bewijzen dat een
waarde correct was en konden verschillende schermen dezelfde berekening anders uitvoeren. In de
huidige architectuur leveren modules zoals \code{lapAnalytics.js}, \code{classBattle.js},
\code{lapPrediction.js} en \code{pitstopPlanner.js} volledige resultaten. De renderer beperkt zich
tot presentatie en gebruikersinteractie.

\section{Vensterarchitectuur}
Er is één hoofdvenster voor de primaire wagen. Wanneer de gebruiker via de setup een tweede of derde
wagen toevoegt, opent voor elke extra wagen een dashboardvenster met dezelfde renderer en een andere
wagenparameter. Alle vensters ontvangen dezelfde collectorstatus. De timingwebsite wordt dus slechts
eenmaal geladen en bevraagd.

Grafieken worden in afzonderlijke vensters geopend. Ook zij gebruiken de reeds opgeslagen lap
history en starten geen eigen collector. Deze aanpak houdt het hoofdscherm compact en maakt het
mogelijk om analysevensters alleen te openen wanneer ze nodig zijn.

\section{Datastroom}
Elke poll doorloopt dezelfde keten:

\begin{enumerate}
  \item Het extractiescript leest tabellen en sessievelden uit de livepagina.
  \item De parser herkent bruikbare headers en zet elke rij om naar canonieke velden.
  \item De storage-normalizer voegt provider- en sessiecontext toe.
  \item De lap history en actuele rijen worden aan de analysemodules doorgegeven.
  \item Voor elke gevolgde wagen worden dashboardanalyse, voorspelling en pitstopplan berekend.
  \item De volledige collectorstatus wordt opgeslagen en via IPC naar alle vensters gestuurd.
\end{enumerate}

De UI en de berekeningen worden dus gevoed door dezelfde consistente snapshot. Dit voorkomt dat een
grafiek bijvoorbeeld met gegevens van een andere poll rekent dan het hoofdscherm.

\begin{figure}[H]
\centering

\begin{tikzpicture}[
  node distance=0.65cm,
  flowbox/.style={
    draw,
    rounded corners=2pt,
    minimum width=9.5cm,
    minimum height=1.05cm,
    align=center,
    fill=blue!5,
    line width=0.8pt
  },
  outputbox/.style={
    draw,
    rounded corners=2pt,
    minimum width=4.4cm,
    minimum height=1.05cm,
    align=center,
    fill=green!8,
    line width=0.8pt
  },
  arrow/.style={
    -{Latex[length=2.5mm]},
    line width=0.8pt
  }
]

\node[flowbox] (source) {
  \textbf{Live timingwebsite}\\
  RIS Timing of GetRaceResults
};

\node[flowbox, below=of source] (parser) {
  \textbf{Providerspecifieke parser}\\
  Leest de DOM-structuur en herkent kolommen en sessiegegevens
};

\node[flowbox, below=of parser] (normalisation) {
  \textbf{Normalisatie}\\
  Zet providerspecifieke waarden om naar één uniform gegevensmodel
};

\node[flowbox, below=of normalisation] (storage) {
  \textbf{Lokale opslag}\\
  Bewaart snapshots, rondehistoriek, metadata en berekende toestand
};

\node[flowbox, below=of storage] (analytics) {
  \textbf{Analytics en domeinlogica}\\
  Geldigheid, gemiddelden, voorspellingen, gaps en pitstopplanning
};

\node[
  outputbox,
  below=0.8cm of analytics,
  xshift=-4cm
] (dashboard) {
  \textbf{Dashboard en grafieken}\\
  Live informatie voor de race-engineer
};

\node[
  outputbox,
  below=0.8cm of analytics,
  xshift=4cm
] (reports) {
  \textbf{PDF-rapporten}\\
  Analyse per stint en volledige sessie
};

\draw[arrow] (source) -- (parser);
\draw[arrow] (parser) -- (normalisation);
\draw[arrow] (normalisation) -- (storage);
\draw[arrow] (storage) -- (analytics);
\draw[arrow] (analytics) -- (dashboard);
\draw[arrow] (analytics) -- (reports);

\end{tikzpicture}

\caption{Datastroom van de live timingpagina naar het dashboard en de rapportering (verwerken data elk op eigen manier).}
\label{fig:datastroom}
\end{figure}

%% ZET DIT IN DE TEXT
%% Figuur~\ref{fig:datastroom} toont hoe de ruwe timingdata stapsgewijs
%% wordt omgezet naar informatie voor het dashboard en de PDF-rapporten.

\section{Instellingen en configuratie}
Instellingen worden als JSON in de Electron user-datafolder bewaard. Ze omvatten timing-URL,
gevolgde wagens, modus, opslagmap, referentietijden en pitstopregels. Referentietijden zijn
afzonderlijk per sessiemodus opgeslagen, zodat een qualifyingreferentie een racereferentie niet
overschrijft. Circuitprofielen voor Spa, Zolder en Assen bevatten standaardwaarden voor de relevante
pitlengte en FCY-snelheid. In de pitstopsetup kunnen deze waarden per evenement worden overschreven
wanneer het reglement afwijkt.
